
\subsection*{Coordinate polari}

Calcoliamo
\[
\iint_{B_R(0)} 1 \, dx\,dy
\]
dove $B_R(0)$ è il disco di raggio $R$ centrato nell'origine.


Cambio di variabili:
\[
\begin{cases}
x = \rho \cos\theta \\
y = \rho \sin\theta
\end{cases}
\]

Dominio:
\[
0 \le \rho \le R,
\qquad
0 \le \theta \le 2\pi
\]

Jacobiano:
\[
J_\varphi =
\begin{vmatrix}
\cos\theta & -\rho\sin\theta \\
\sin\theta & \rho\cos\theta
\end{vmatrix}
= \rho
\]

Quindi:
\[
dx\,dy = \rho \, d\rho\, d\theta
\]

L’integrale diventa:
\[
\iint_{B_R(0)} 1\, dx\,dy
=
\int_0^{2\pi} \int_0^R \rho \, d\rho\, d\theta
\]

Calcolo:
\[
\int_0^R \rho \, d\rho = \frac{R^2}{2}
\]

\[
\int_0^{2\pi} d\theta = 2\pi
\]

Risultato:
\[
\iint_{B_R(0)} 1\, dx\,dy
=
\frac{R^2}{2} \cdot 2\pi
=
\pi R^2
\]

\section*{Coordinate cilindriche}

\[
\begin{cases}
x = \rho \cos\theta \\
y = \rho \sin\theta \\
z = z
\end{cases}
\]

Jacobiano:
\[
J = \rho
\]

Elemento di volume:
\[
dx\,dy\,dz = \rho \, d\rho\, d\theta\, dz
\]

\section*{Coordinate sferiche}

%\documentclass[tikz,border=5pt]{standalone}
%\usepackage{amsmath}
%\usetikzlibrary{arrows.meta}

\begin{figure}
\begin{tikzpicture}[scale=3, >=Latex]
% Assi
\draw[->] (0,0,0) -- (1.2,0,0) node[below] {$x$};
\draw[->] (0,0,0) -- (0,1.2,0) node[left] {$y$};
\draw[->] (0,0,0) -- (0,0,1.2) node[left] {$z$};

% Parametri del punto
\def\r{1}
\def\theta{40}   % angolo azimutale
\def\phi{50}     % angolo polare

% Coordinate cartesiane del punto
\pgfmathsetmacro{\px}{\r*sin(\phi)*cos(\theta)}
\pgfmathsetmacro{\py}{\r*sin(\phi)*sin(\theta)}
\pgfmathsetmacro{\pz}{\r*cos(\phi)}

% Punto P
\coordinate (P) at (\px,\py,\pz);

% Raggio rho
\draw[thick] (0,0,0) -- (P) node[midway, right] {$\rho$};

% Proiezione sul piano xy
\coordinate (Q) at (\px,\py,0);
\draw[dashed] (P) -- (Q);
\draw[dashed] (0,0,0) -- (Q);

% Angolo theta nel piano xy
\draw (0.3,0,0) arc[start angle=0, end angle=\theta, radius=0.3];
\node at (0.45,0.15,0) {$\theta$};

% Angolo phi dal semiasse z
\draw (0,0,0.4) arc[start angle=90, end angle={90-\phi}, radius=0.4];
\node at (0.15,0,0.55) {$\varphi$};

% Punto P
\fill (P) circle (0.5pt);
\node[right] at (P) {$P(\rho,\theta,\varphi)$};

\end{tikzpicture}
\end{figure}

\[
\begin{cases}
x = \rho \sin\varphi \cos\theta \\
y = \rho \sin\varphi \sin\theta \\
z = \rho \cos\varphi
\end{cases}
\]

con:
\[
\rho \ge 0,
\qquad
0 \le \theta \le 2\pi,
\qquad
0 \le \varphi \le \pi
\]

Jacobiano:
\[
J = \rho^2 \sin\varphi
\]

Elemento di volume:
\[
dx\,dy\,dz
=
\rho^2 \sin\varphi
\, d\rho\, d\varphi\, d\theta
\]

\section*{Formula generale di cambio di variabili}

Se $\varphi : U \to \mathbb{R}^n$ è un diffeomorfismo regolare, allora:

\[
\int_{\varphi(U)}
f(x)\, dx
=
\int_U
f(\varphi(u))
\left| \det D\varphi(u) \right|
\, du
\]
